\section{Sequential Pattern Mining (PrefixSpan)}
\label{sec:prefixspan}

Clusters answer the question \emph{which patients look alike}. They say nothing about \emph{the order in which clinical events unfold}. The second analytical module addresses precisely that, mining frequent ordered sub-sequences from the patient timelines built in Section~5.2.3 with \textbf{PrefixSpan}~\cite{pei2004}. The implementation is a single script, \texttt{prefix\_span.py} (496 lines).

\subsection{Mining Pipeline Overview}

\begin{figure}[H]
\centering
\includegraphics[width=0.95\linewidth]{figs/prefix_span_pipeline.png}
\caption{Phases of the sequence-mining pipeline.}
\label{fig:prefixspan-flow}
\end{figure}

\textbf{Integer encoding.} PrefixSpan compares items by identity; raw strings would be both slow and memory-heavy. A single bidirectional dictionary maps every unique \texttt{C:/M:/P:/CP:} item to a small integer; patient timelines are encoded to lists of integers, and the inverse map decodes mined patterns back to human-readable strings.

\textbf{Top-$k$ closed mining.}
\begin{verbatim}
PrefixSpan(seqs).topk(
    k=TOP_K * OVERMINE_MULTIPLIER,   # 500 * 20 = 10000
    closed=True,
)
\end{verbatim}
\texttt{closed=True} discards any pattern that is the strict prefix of another with the same support, eliminating a large class of trivial duplicates. The \texttt{OVERMINE\_MULTIPLIER} is the key idea: the algorithm is asked for twenty times more patterns than will ever be shown, so the post-mining filters of Sections~\ref{subsec:gml} and~\ref{subsec:dedup} have a healthy candidate pool.

\textbf{Hard filters.} Before reranking:
\begin{itemize}
    \item \textbf{Support}: $\texttt{support\_pct} \geq 0.01$ and absolute support $\geq 2$.
    \item \textbf{Length}: $3 \leq |p| \leq 5$. Shorter is too generic, longer too rare.
    \item \textbf{Type cadence}: any pattern with three consecutive same-type items is rejected (e.g.\ three conditions in a row), since these almost always reflect a single billing episode.
\end{itemize}
